% ============================================================
%  PriceIQ — Software Requirements Specification
%  BTech Final Year Project, 2025-2026
% ============================================================
\documentclass[12pt,a4paper]{report}

\usepackage[top=1in,bottom=1in,left=1.2in,right=1in]{geometry}
\usepackage{booktabs}
\usepackage{longtable}
\usepackage{array}
\usepackage{tabularx}
\usepackage{multirow}
\usepackage[hidelinks]{hyperref}
\usepackage{xcolor}
\usepackage{enumitem}
\usepackage{fancyhdr}
\usepackage{parskip}
\usepackage{titlesec}
\usepackage{setspace}
\usepackage{mdframed}

\definecolor{priblue}{RGB}{0,82,165}
\definecolor{rowgray}{RGB}{245,245,245}
\definecolor{headerblue}{RGB}{220,235,255}

\pagestyle{fancy}
\fancyhf{}
\fancyhead[L]{\small\textit{PriceIQ --- Software Requirements Specification v1.0}}
\fancyhead[R]{\small\textit{BTech Final Year Project 2025--2026}}
\fancyfoot[C]{\thepage}
\renewcommand{\headrulewidth}{0.4pt}

\setlength{\parskip}{6pt}
\setlength{\parindent}{0pt}
\onehalfspacing

\newcolumntype{L}[1]{>{\raggedright\arraybackslash}p{#1}}
\newcolumntype{C}[1]{>{\centering\arraybackslash}p{#1}}

% Requirement box
\newmdenv[linecolor=priblue,backgroundcolor=headerblue!30,
          linewidth=1pt,leftmargin=0,rightmargin=0,
          skipabove=6pt,skipbelow=6pt]{reqbox}

\begin{document}

% ─── TITLE PAGE ─────────────────────────────────────────────────────────────
\begin{titlepage}
  \centering
  \vspace*{1.8cm}
  {\Huge\bfseries\color{priblue} PriceIQ\par}
  \vspace{0.4cm}
  {\large AI-Powered Electronics Price Comparison for Indian Consumers\par}
  \vspace{1.4cm}
  \rule{0.85\linewidth}{1.5pt}\\[0.5cm]
  {\LARGE\bfseries Software Requirements Specification (SRS)\par}
  \vspace{0.3cm}
  {\large Document Version 1.0\par}
  \rule{0.85\linewidth}{1.5pt}
  \vspace{1.8cm}

  \begin{tabular}{@{}ll@{}}
    \textbf{Project Type:}   & BTech Final Year Project \\[4pt]
    \textbf{Department:}     & Computer Science and Engineering \\[4pt]
    \textbf{Academic Year:}  & 2025--2026 \\[4pt]
    \textbf{Prepared By:}    & Adrij Bhadra \\[4pt]
    \textbf{Date:}           & April 2026 \\[4pt]
    \textbf{Document Status:}& Final \\
  \end{tabular}

  \vfill
  {\small This document is prepared for academic evaluation only.\\
  PriceIQ is an independently developed final-year academic project.}
\end{titlepage}

\tableofcontents
\newpage

% ─── CHAPTER 1: INTRODUCTION ─────────────────────────────────────────────────
\chapter{Introduction}

\section{Purpose}
This Software Requirements Specification (SRS) document formally describes the
functional and non-functional requirements for \textbf{PriceIQ}, an AI-powered
electronics price comparison platform designed for Indian consumers. The document
is intended for the following audiences:

\begin{itemize}[noitemsep]
  \item Academic evaluators and examiners assessing the BTech final-year project.
  \item The developer(s) implementing and maintaining the system.
  \item Any future contributors who extend the platform.
\end{itemize}

This SRS is the authoritative reference for what the system \emph{must} do, how
it must behave, and the constraints under which it operates. It does not describe
implementation details, which are covered in the System Design Document.

\section{Scope}
PriceIQ is a web-based application that allows users to search for consumer
electronics and compare real-time prices from three major Indian online retailers:
Amazon India (\texttt{amazon.in}), Flipkart (\texttt{flipkart.com}), and Vijay
Sales (\texttt{vijaysales.com}). The system additionally generates AI-driven
purchase recommendations using the Groq LLaMA-3 large language model.

The application is scoped to the Indian market (prices in Indian Rupees, INR),
and covers the following product categories: Smartphones, Laptops, Headphones,
Televisions, Tablets, Cameras, and Wearables.

\textbf{In scope:}
\begin{itemize}[noitemsep]
  \item User search and live scraping across three platforms.
  \item Price comparison display with source-level detail.
  \item Price history tracking across scrape cycles.
  \item AI-generated product summary and purchase recommendation.
  \item Batch scraping of a curated product catalogue.
  \item Administrative dashboard with KPI metrics.
\end{itemize}

\textbf{Out of scope:}
\begin{itemize}[noitemsep]
  \item E-commerce transaction or cart functionality.
  \item User account management or personalisation.
  \item Price alert notifications or email/SMS delivery.
  \item Support for non-Indian retailers or non-INR currencies.
  \item Mobile (iOS/Android) native applications.
\end{itemize}

\section{Definitions, Acronyms, and Abbreviations}

\begin{longtable}{L{2.8cm} L{10cm}}
  \toprule
  \textbf{Term} & \textbf{Definition} \\
  \midrule
  \endhead
  PriceIQ & The name of the system being specified. \\
  SRS & Software Requirements Specification. \\
  API & Application Programming Interface. \\
  LLM & Large Language Model. \\
  Groq & A hardware-accelerated AI inference provider; used to run LLaMA-3. \\
  LLaMA-3 & Open-weight large language model developed by Meta. \\
  ORM & Object-Relational Mapper; specifically Prisma 7 in this project. \\
  SSR & Server-Side Rendering (Next.js feature). \\
  ASIN & Amazon Standard Identification Number, a unique product identifier on Amazon. \\
  Playwright & A browser automation library used to scrape Flipkart. \\
  Slug & A URL-safe lowercase string uniquely identifying a product, e.g., \texttt{iphone-15-128gb}. \\
  INR & Indian Rupee, the currency used throughout the system. \\
  FR & Functional Requirement. \\
  NFR & Non-Functional Requirement. \\
  MRP & Maximum Retail Price, the government-mandated ceiling price on Indian goods. \\
  \bottomrule
\end{longtable}

\section{Overview of Document}
The remainder of this document is organised as follows:

\begin{itemize}[noitemsep]
  \item \textbf{Chapter 2} provides an overall description of the product,
        including product perspective, key functions, user classes, and
        operating environment.
  \item \textbf{Chapter 3} specifies all functional requirements, grouped
        into ten major feature areas with detailed use-case descriptions.
  \item \textbf{Chapter 4} defines non-functional requirements covering
        performance, reliability, security, usability, and maintainability.
  \item \textbf{Chapter 5} describes external interface requirements including
        the user interface, third-party API interfaces, and communication
        protocols.
  \item \textbf{Chapter 6} presents the system's data model and persistence
        requirements.
  \item \textbf{Chapter 7} lists assumptions, dependencies, and known
        constraints.
\end{itemize}

% ─── CHAPTER 2: OVERALL DESCRIPTION ─────────────────────────────────────────
\chapter{Overall Description}

\section{Product Perspective}
PriceIQ is a standalone web application operating in a client-server model. It
is not a sub-system of any larger platform; however, it depends on three external
data sources (Amazon India, Flipkart, Vijay Sales) and one external AI inference
service (Groq API) to deliver its core value.

The system fills a gap in the Indian consumer market: existing price-comparison
services (e.g., BuyHatke, PriceDekho) do not provide AI-generated analysis
tailored to Indian buying priorities such as warranty support, seller trust,
and delivery speed. PriceIQ combines real-time scraping with LLM inference to
address this gap.

\begin{figure}[h]
\centering
\textit{(Refer to the System Design Document for the full architecture diagram.)}
\end{figure}

\section{Product Functions}
The major functions provided by PriceIQ are summarised below:

\begin{enumerate}[noitemsep]
  \item \textbf{Live Search:} On receiving a user query, the system concurrently
        scrapes Amazon India, Flipkart, and Vijay Sales and returns normalised
        price data.
  \item \textbf{Product Identity Matching:} A scoring engine parses the query
        into brand, model, storage, and variant tokens and scores each retailer
        candidate to prevent wrong-variant or accessory listings.
  \item \textbf{Price Comparison:} Aggregated listings are displayed in a
        structured comparison table showing price, rating, stock status, and
        a ``best deal'' highlight.
  \item \textbf{AI Recommendation:} The Groq LLaMA-3 model analyses the
        available listings and generates a concise purchase recommendation
        tailored to Indian consumers.
  \item \textbf{Price History:} Every scrape cycle appends a timestamped price
        record, enabling users to observe trends via a line chart.
  \item \textbf{Batch Scraping:} An administrative endpoint runs a full scrape
        of a curated product catalogue and populates the database.
  \item \textbf{Search:} Users can query the product database with text search
        and apply category, source, and sort filters.
  \item \textbf{Dashboard:} Administrators can view system KPIs including total
        products, listing counts by source, recent searches, and
        category distribution.
\end{enumerate}

\section{User Classes and Characteristics}

\begin{longtable}{L{2.5cm} L{3cm} L{7cm}}
  \toprule
  \textbf{User Class} & \textbf{Frequency} & \textbf{Characteristics} \\
  \midrule
  \endhead
  General Consumer &
  High &
  Indian consumer looking to compare prices before purchasing electronics.
  Non-technical. Interacts via search bar and product page. \\[4pt]
  Power User &
  Medium &
  Tech-savvy user who reads AI analysis, checks price history charts, and
  uses category/sort filters. \\[4pt]
  Administrator &
  Low &
  Developer or project maintainer who triggers batch scraping, monitors
  the dashboard, and manages the product catalogue. \\
  \bottomrule
\end{longtable}

\section{Operating Environment}
The system operates in the following environment:

\begin{itemize}[noitemsep]
  \item \textbf{Server OS:} Linux (Docker container) or Windows 11 (development).
  \item \textbf{Runtime:} Node.js 20+, Next.js 16 App Router.
  \item \textbf{Database:} PostgreSQL 16 via Docker Compose.
  \item \textbf{Browser:} Any modern Chromium, Firefox, or Safari-based browser
        (desktop and mobile).
  \item \textbf{Network:} Outbound HTTPS access to \texttt{amazon.in},
        \texttt{flipkart.com}, \texttt{vsprod.vijaysales.com}, and
        \texttt{api.groq.com}.
\end{itemize}

\section{Design and Implementation Constraints}
\begin{itemize}[noitemsep]
  \item The system must use Next.js 16 App Router architecture to comply with
        the project's approved technology stack.
  \item All monetary values are stored and displayed in INR only.
  \item Flipkart scraping requires a headless Chromium browser (Playwright)
        and cannot be replaced by a plain HTTP client due to the site's
        bot-detection measures.
  \item The Groq API key must be stored as an environment variable and never
        committed to version control.
  \item The database schema is managed by Prisma 7; the \texttt{DATABASE\_URL}
        must be declared in \texttt{prisma.config.ts} (not \texttt{schema.prisma}).
  \item The batch scrape endpoint must be protected by a shared secret
        (\texttt{SCRAPE\_SECRET}) transmitted via the \texttt{x-scrape-secret}
        HTTP header.
\end{itemize}

% ─── CHAPTER 3: FUNCTIONAL REQUIREMENTS ─────────────────────────────────────
\chapter{Functional Requirements}

Each requirement is given a unique identifier in the format \textbf{FR-NNN}.
Priority levels: \textbf{M}~=~Must Have, \textbf{S}~=~Should Have,
\textbf{C}~=~Could Have.

% ── FR-001 Live Search ──
\section{FR-001: Live Search}

\begin{reqbox}
\textbf{ID:} FR-001 \quad \textbf{Priority:} M \quad \textbf{Actor:} General Consumer
\end{reqbox}

\textbf{Description:} When a user submits a search query that returns no results
from the local database, the system must automatically initiate a live search
by concurrently scraping Amazon India, Flipkart, and Vijay Sales.

\begin{longtable}{L{3cm} L{10cm}}
  \toprule
  \textbf{Field} & \textbf{Detail} \\
  \midrule
  \endhead
  Pre-conditions & Application is running; user has entered a non-empty search
                   query; no matching product exists in the database. \\[3pt]
  Main Flow &
    1. User submits query via the search bar. \newline
    2. System queries the database; result count is zero. \newline
    3. \texttt{LiveSearchTrigger} component fires a POST to
       \texttt{/api/search/live}. \newline
    4. Server calls \texttt{liveSearch(query)}: Vijay Sales (GraphQL) and
       Amazon (HTTP) are scraped in parallel; Flipkart (Playwright) runs
       sequentially after. \newline
    5. Each platform returns up to one \texttt{NormalizedProduct}. \newline
    6. System upserts a canonical \texttt{Product} record and one
       \texttt{ProductListing} per platform. \newline
    7. A \texttt{PriceHistory} entry is recorded for each listing. \newline
    8. System redirects the user to \texttt{/product/\{id\}}. \\[3pt]
  Alternative Flow A &
    Platform scrape fails (timeout, bot block): that platform is skipped;
    remaining platforms' results are still returned. \\[3pt]
  Alternative Flow B &
    No platform returns an accepted candidate: system returns an error message
    to the user suggesting they refine the query. \\[3pt]
  Post-conditions & At least one \texttt{ProductListing} record is stored; user
                    is shown the product comparison page. \\
  \bottomrule
\end{longtable}

% ── FR-002 Product Identity Matching ──
\section{FR-002: Product Identity Matching}

\begin{reqbox}
\textbf{ID:} FR-002 \quad \textbf{Priority:} M \quad \textbf{Actor:} System (internal)
\end{reqbox}

\textbf{Description:} The system must use a structured scoring engine to
verify that a retailer-returned product title matches the user's search intent
before storing or displaying the result.

\begin{longtable}{L{3cm} L{10cm}}
  \toprule
  \textbf{Field} & \textbf{Detail} \\
  \midrule
  \endhead
  Pre-conditions & A raw candidate product title and a parsed query object are
                   available. \\[3pt]
  Main Flow &
    1. \texttt{parseQuery()} extracts brand, model tokens, storage token (e.g.,
       128GB), and variant token (e.g., Ultra) from the query. \newline
    2. \texttt{scoreCandidate()} evaluates the candidate title: \newline
       \quad a) Hard-reject accessory keywords (case, cover, charger, etc.). \newline
       \quad b) Hard-reject renewed/refurbished listings unless query requests
              them. \newline
       \quad c) Hard-reject if no model tokens match. \newline
       \quad d) Hard-reject storage or variant conflicts. \newline
       \quad e) Apply soft score bonuses for brand match (+30), model token
              coverage (+40 max), storage match (+20), variant match (+20). \newline
       \quad f) Apply soft penalties for combos, unexpected variants, missing
              storage, missing brand. \newline
    3. A candidate is accepted if score $\geq 40$ (when model tokens exist) or
       score $\geq 20$ (vague query). \newline
    4. The highest-scoring accepted candidate is selected; ties are broken by
       review count. \\[3pt]
  Post-conditions & Only products that semantically match the query are stored
                    in the database. \\
  \bottomrule
\end{longtable}

% ── FR-003 Price Comparison Display ──
\section{FR-003: Price Comparison Display}

\begin{reqbox}
\textbf{ID:} FR-003 \quad \textbf{Priority:} M \quad \textbf{Actor:} General Consumer
\end{reqbox}

\textbf{Description:} The product detail page must display a structured
comparison of all available platform listings for a product.

\begin{longtable}{L{3cm} L{10cm}}
  \toprule
  \textbf{Field} & \textbf{Detail} \\
  \midrule
  \endhead
  Pre-conditions & A product with at least one listing exists in the database. \\[3pt]
  Main Flow &
    1. User navigates to \texttt{/product/\{id\}}. \newline
    2. System loads \texttt{Product} with all \texttt{ProductListing} records. \newline
    3. Listings are sorted: in-stock first, then by ascending price. \newline
    4. The cheapest in-stock listing is highlighted as ``Best Deal''. \newline
    5. Price differences from the best deal are shown as percentage. \newline
    6. Each listing shows: platform name, price (INR), star rating,
       review count, MRP (Vijay Sales only), stock status, and an
       outbound link. \\[3pt]
  Alternative Flow &
    A listing is out of stock: it is shown last with a grey ``Out of Stock'' badge. \\[3pt]
  Post-conditions & User can see prices from all scraped platforms on one page. \\
  \bottomrule
\end{longtable}

% ── FR-004 AI Purchase Recommendation ──
\section{FR-004: AI Purchase Recommendation}

\begin{reqbox}
\textbf{ID:} FR-004 \quad \textbf{Priority:} M \quad \textbf{Actor:} General Consumer, System
\end{reqbox}

\textbf{Description:} The system must generate an AI-driven purchase
recommendation for each product using the Groq LLaMA-3 inference API.

\begin{longtable}{L{3cm} L{10cm}}
  \toprule
  \textbf{Field} & \textbf{Detail} \\
  \midrule
  \endhead
  Pre-conditions & Product has at least one listing; Groq API key is configured. \\[3pt]
  Main Flow &
    1. After live search completes, server calls \texttt{analyzeProduct()} with
       product name, category, domain, and all listings. \newline
    2. Groq LLaMA-3.3-70b model receives a structured prompt requesting JSON
       output with: \texttt{summary}, \texttt{recommendation},
       \texttt{bestSource}, \texttt{confidenceNote}. \newline
    3. Response is parsed; if JSON is malformed, a regex-based text fallback
       extracts the fields. \newline
    4. Result is stored in \texttt{AIAnalysis}; cached for 24~hours. \newline
    5. User sees the analysis on the ``AI Analysis'' tab of the product page. \\[3pt]
  Alternative Flow &
    Groq API call fails: product is still saved; AI tab shows a ``Refresh''
    button so the user can retry. \\[3pt]
  Post-conditions & An \texttt{AIAnalysis} record is stored with a 24-hour TTL. \\
  \bottomrule
\end{longtable}

% ── FR-005 Price History ──
\section{FR-005: Price History Tracking}

\begin{reqbox}
\textbf{ID:} FR-005 \quad \textbf{Priority:} S \quad \textbf{Actor:} System, Power User
\end{reqbox}

\textbf{Description:} Every scrape event must append a timestamped price record
to the \texttt{PriceHistory} table, enabling trend visualisation.

\begin{longtable}{L{3cm} L{10cm}}
  \toprule
  \textbf{Field} & \textbf{Detail} \\
  \midrule
  \endhead
  Pre-conditions & A scrape event (live or batch) has completed with at least
                   one accepted listing. \\[3pt]
  Main Flow &
    1. For each accepted \texttt{NormalizedProduct}, system creates a
       \texttt{PriceHistory} row with: \texttt{productId}, \texttt{source},
       \texttt{price}, \texttt{recordedAt} (UTC timestamp). \newline
    2. Rows are never updated or deleted (append-only). \newline
    3. Product page renders a line chart via Recharts querying
       \texttt{/api/price-history/\{id\}}. \\[3pt]
  Post-conditions & Chart on product page reflects all historical price points. \\
  \bottomrule
\end{longtable}

% ── FR-006 Batch Scrape ──
\section{FR-006: Batch Scraping}

\begin{reqbox}
\textbf{ID:} FR-006 \quad \textbf{Priority:} M \quad \textbf{Actor:} Administrator
\end{reqbox}

\textbf{Description:} An authenticated HTTP endpoint must trigger a full scrape
of the curated product catalogue (\texttt{products.ts}) across all three
platforms.

\begin{longtable}{L{3cm} L{10cm}}
  \toprule
  \textbf{Field} & \textbf{Detail} \\
  \midrule
  \endhead
  Pre-conditions & Request includes a valid \texttt{x-scrape-secret} header. \\[3pt]
  Main Flow &
    1. POST to \texttt{/api/scrape} with header
       \texttt{x-scrape-secret: <SCRAPE\_SECRET>}. \newline
    2. Server validates the secret against \texttt{process.env.SCRAPE\_SECRET}. \newline
    3. \texttt{runAllScrapers()} is called: iterates over 4 products $\times$ 3
       sources. \newline
    4. Results are grouped by canonical slug, upserted into \texttt{Product}
       and \texttt{ProductListing}, and price history is recorded. \newline
    5. JSON response returns the count of upserted records. \\[3pt]
  Alternative Flow &
    Missing or wrong secret: server responds 403 Forbidden immediately. \\[3pt]
  Post-conditions & Database reflects latest prices for all catalogue products. \\
  \bottomrule
\end{longtable}

% ── FR-007 Database Search ──
\section{FR-007: Database Text Search and Filtering}

\begin{reqbox}
\textbf{ID:} FR-007 \quad \textbf{Priority:} M \quad \textbf{Actor:} General Consumer
\end{reqbox}

\textbf{Description:} The search results page must support full-text search of
the product database with category filtering, source filtering, and multiple
sort options.

\begin{longtable}{L{3cm} L{10cm}}
  \toprule
  \textbf{Field} & \textbf{Detail} \\
  \midrule
  \endhead
  Pre-conditions & Database contains at least one product. \\[3pt]
  Main Flow &
    1. User enters query in the search bar and optionally selects filters. \newline
    2. GET \texttt{/api/search?q=\{query\}\&category=\{cat\}\&sort=\{sort\}} is issued. \newline
    3. Database performs case-insensitive match on \texttt{name}, \texttt{brand},
       \texttt{category}, and \texttt{description} fields. \newline
    4. Optional filters: category (exact match), source (listing source). \newline
    5. Sort options: price ascending/descending, rating, source count (default). \newline
    6. Search query is logged to \texttt{SearchLog}. \newline
    7. Results rendered as a grid of \texttt{ProductCard} components. \\[3pt]
  Alternative Flow &
    Zero results found: \texttt{LiveSearchTrigger} is shown (see FR-001). \\[3pt]
  Post-conditions & Matching products displayed; query logged. \\
  \bottomrule
\end{longtable}

% ── FR-008 Refresh Prices ──
\section{FR-008: On-Demand Price Refresh}

\begin{reqbox}
\textbf{ID:} FR-008 \quad \textbf{Priority:} S \quad \textbf{Actor:} Power User
\end{reqbox}

\textbf{Description:} From the product detail page, a user must be able to
trigger a re-scrape of all three platforms for the displayed product.

\begin{longtable}{L{3cm} L{10cm}}
  \toprule
  \textbf{Field} & \textbf{Detail} \\
  \midrule
  \endhead
  Pre-conditions & Product page is open with an existing product record. \\[3pt]
  Main Flow &
    1. User clicks ``Refresh Prices'' button. \newline
    2. Frontend calls \texttt{/api/product/\{id\}/refresh}. \newline
    3. Server re-runs \texttt{liveSearch()} with the product's canonical name. \newline
    4. Updated listings are upserted; new price history rows are created. \newline
    5. Page re-renders with updated data; toast notification confirms success. \\[3pt]
  Post-conditions & Listings and price history updated; AI analysis cache invalidated. \\
  \bottomrule
\end{longtable}

% ── FR-009 Dashboard ──
\section{FR-009: Administrative Dashboard}

\begin{reqbox}
\textbf{ID:} FR-009 \quad \textbf{Priority:} S \quad \textbf{Actor:} Administrator
\end{reqbox}

\textbf{Description:} A dedicated dashboard page must display system-wide KPI
metrics including product count, listing distribution by source, recent search
activity, and category breakdown.

\begin{longtable}{L{3cm} L{10cm}}
  \toprule
  \textbf{Field} & \textbf{Detail} \\
  \midrule
  \endhead
  Pre-conditions & At least one product and one search log entry exist. \\[3pt]
  Main Flow &
    1. User navigates to \texttt{/dashboard}. \newline
    2. Page calls \texttt{/api/dashboard} which aggregates: total products,
       total listings per source, recent search queries, and category
       distribution. \newline
    3. KPI cards, a source distribution pie chart (Recharts), a category
       bar chart, and a recent-searches table are rendered. \\[3pt]
  Post-conditions & Administrator can assess system coverage at a glance. \\
  \bottomrule
\end{longtable}

% ── FR-010 Dark Mode ──
\section{FR-010: Dark Mode UI Theme}

\begin{reqbox}
\textbf{ID:} FR-010 \quad \textbf{Priority:} C \quad \textbf{Actor:} General Consumer
\end{reqbox}

\textbf{Description:} The application UI must support a dark colour theme that
is applied by default and persists across sessions.

\begin{longtable}{L{3cm} L{10cm}}
  \toprule
  \textbf{Field} & \textbf{Detail} \\
  \midrule
  \endhead
  Pre-conditions & User accesses any page. \\[3pt]
  Main Flow &
    1. \texttt{ThemeProvider} (next-themes) initialises with theme = ``dark''. \newline
    2. Theme is stored in \texttt{localStorage} and restored on subsequent visits. \newline
    3. All Tailwind CSS and Radix UI components respond to the \texttt{dark:}
       variant class. \\[3pt]
  Post-conditions & UI renders in dark mode consistently across all pages. \\
  \bottomrule
\end{longtable}

% ─── CHAPTER 4: NON-FUNCTIONAL REQUIREMENTS ──────────────────────────────────
\chapter{Non-Functional Requirements}

\section{Performance Requirements}

\begin{longtable}{L{2cm} L{4cm} L{7cm}}
  \toprule
  \textbf{ID} & \textbf{Requirement} & \textbf{Metric / Acceptance Criterion} \\
  \midrule
  \endhead
  NFR-P01 & Live search latency &
    Total live search (all 3 platforms) must complete within \textbf{60 seconds}
    under normal network conditions. Flipkart (Playwright) dominates latency. \\[3pt]
  NFR-P02 & Cached search latency &
    Database text search must return results within \textbf{500 ms} for a
    product catalogue of up to 1\,000 products. \\[3pt]
  NFR-P03 & Product page load &
    SSR product page must reach Time-to-First-Byte within \textbf{2 seconds}
    on a standard broadband connection. \\[3pt]
  NFR-P04 & AI analysis latency &
    Groq API call must complete within \textbf{8 seconds} at the model
    \texttt{llama-3.3-70b-versatile} with a 450-token output limit. \\[3pt]
  NFR-P05 & Concurrent users &
    System must handle at least \textbf{5 concurrent live search requests}
    without returning server errors, given the single-instance deployment model. \\
  \bottomrule
\end{longtable}

\section{Reliability Requirements}

\begin{longtable}{L{2cm} L{4cm} L{7cm}}
  \toprule
  \textbf{ID} & \textbf{Requirement} & \textbf{Metric / Acceptance Criterion} \\
  \midrule
  \endhead
  NFR-R01 & Partial scrape failure &
    If one or two platform scrapers fail, the remaining results must still be
    saved and shown to the user (graceful degradation). \\[3pt]
  NFR-R02 & AI failure isolation &
    If the Groq API is unreachable, the product page must still render
    price data; the AI tab shows a retry button. \\[3pt]
  NFR-R03 & Database uptime &
    The PostgreSQL instance (Docker) must recover automatically on container
    restart. Data must survive container restarts via Docker volume. \\[3pt]
  NFR-R04 & Bot detection handling &
    If Amazon.in returns a 503 or CAPTCHA, the Amazon scraper must log a
    warning and return null rather than throwing an unhandled exception. \\
  \bottomrule
\end{longtable}

\section{Security Requirements}

\begin{longtable}{L{2cm} L{4cm} L{7cm}}
  \toprule
  \textbf{ID} & \textbf{Requirement} & \textbf{Metric / Acceptance Criterion} \\
  \midrule
  \endhead
  NFR-S01 & Scrape endpoint protection &
    POST \texttt{/api/scrape} must reject any request whose
    \texttt{x-scrape-secret} header does not match
    \texttt{process.env.SCRAPE\_SECRET}. Response: 403 Forbidden. \\[3pt]
  NFR-S02 & API key security &
    \texttt{GROQ\_API\_KEY}, \texttt{DATABASE\_URL}, and \texttt{SCRAPE\_SECRET}
    must be stored only in \texttt{.env.local} (or Docker environment variables)
    and must never be embedded in client-side code. \\[3pt]
  NFR-S03 & Input sanitisation &
    Search query parameters must be passed to the Prisma ORM as parameterised
    values, preventing SQL injection. \\[3pt]
  NFR-S04 & No sensitive data in responses &
    API responses must never expose environment variable values,
    database connection strings, or internal stack traces to the client. \\
  \bottomrule
\end{longtable}

\section{Usability Requirements}

\begin{longtable}{L{2cm} L{4cm} L{7cm}}
  \toprule
  \textbf{ID} & \textbf{Requirement} & \textbf{Metric / Acceptance Criterion} \\
  \midrule
  \endhead
  NFR-U01 & Mobile responsiveness &
    All pages must be fully usable on a 375~px wide viewport (iPhone SE)
    using Tailwind CSS responsive utilities. \\[3pt]
  NFR-U02 & Loading feedback &
    The \texttt{LiveSearchTrigger} component must display an animated
    per-platform progress indicator while scraping is in progress. \\[3pt]
  NFR-U03 & Error messaging &
    All user-visible errors must be described in plain English with actionable
    guidance (e.g., ``No results found. Try a more specific query.''). \\[3pt]
  NFR-U04 & Accessibility &
    All interactive controls must have visible focus states and ARIA labels
    sufficient to pass basic keyboard navigation. \\
  \bottomrule
\end{longtable}

\section{Maintainability Requirements}

\begin{longtable}{L{2cm} L{4cm} L{7cm}}
  \toprule
  \textbf{ID} & \textbf{Requirement} & \textbf{Metric / Acceptance Criterion} \\
  \midrule
  \endhead
  NFR-M01 & TypeScript strict mode &
    All source files must compile with \texttt{strict: true} in
    \texttt{tsconfig.json} without type errors. \\[3pt]
  NFR-M02 & Scraper extensibility &
    Adding a new retailer must require changes to at most three files:
    a new scraper module, \texttt{index.ts}, and \texttt{live.ts}. The
    \texttt{matcher.ts} scoring engine and \texttt{types.ts} interface must
    not require modification. \\[3pt]
  NFR-M03 & Product catalogue extension &
    Adding a new product to the curated catalogue must require only adding
    one entry to \texttt{products.ts}; no other code changes. \\
  \bottomrule
\end{longtable}

% ─── CHAPTER 5: EXTERNAL INTERFACE REQUIREMENTS ──────────────────────────────
\chapter{External Interface Requirements}

\section{User Interface Requirements}
\begin{itemize}[noitemsep]
  \item The application must provide a single-page hero section on the home
        route (\texttt{/}) with a prominently placed search bar.
  \item The search bar must support auto-submission on Enter key press and
        display search suggestions fetched from
        \texttt{/api/search/suggestions}.
  \item The product detail page (\texttt{/product/[id]}) must organise content
        into two tabs: \emph{Price Comparison} and \emph{AI Analysis}.
  \item All pages must include a persistent navigation bar with the PriceIQ
        brand logo, a search field, and a link to the dashboard.
  \item Toast notifications (Sonner library) must be used for scrape-trigger
        and refresh-success feedback; no modal dialogs.
\end{itemize}

\section{Hardware Interfaces}
The system has no direct hardware interface requirements beyond standard web
server provisioning. The Playwright headless Chromium browser requires a minimum
of 512~MB RAM per browser instance.

\section{Software Interfaces}

\begin{longtable}{L{3cm} L{3cm} L{7cm}}
  \toprule
  \textbf{Interface} & \textbf{Version} & \textbf{Description} \\
  \midrule
  \endhead
  Groq API &
  \texttt{groq-sdk} 1.1.2 &
  REST-over-HTTPS chat completions API. Model: \texttt{llama-3.3-70b-versatile}.
  Requires \texttt{GROQ\_API\_KEY} in environment. \\[3pt]
  Vijay Sales GraphQL &
  Internal &
  GraphQL endpoint at \texttt{vsprod.vijaysales.com/graphql}. Publicly
  accessible; no auth required. Requires a cache-bypass timestamp
  parameter (\texttt{tp}). \\[3pt]
  Amazon India &
  Public HTML &
  Standard HTTPS fetch of \texttt{amazon.in/s} (search) and
  \texttt{amazon.in/dp/\{ASIN\}} (product page). Subject to bot-detection
  blocking (503 responses). \\[3pt]
  Flipkart &
  Playwright DOM &
  Headless Chromium navigates \texttt{flipkart.com/search}. Anti-detection
  headers and navigator spoof required. \\[3pt]
  Prisma ORM &
  7.7.0 &
  Prisma client with \texttt{@prisma/adapter-pg} (PostgreSQL adapter). \\[3pt]
  PostgreSQL &
  16 &
  Relational database running in Docker. Accessed via connection string. \\[3pt]
  Next.js &
  16.2.4 &
  Full-stack React framework. Provides API routes, SSR, and static assets. \\
  \bottomrule
\end{longtable}

\section{Communication Interfaces}
\begin{itemize}[noitemsep]
  \item All client-server communication uses \textbf{HTTPS} (TLS 1.2+) in
        production; HTTP is acceptable in local development only.
  \item API routes follow \textbf{REST} conventions: POST for mutations,
        GET for queries, standard HTTP status codes.
  \item Scraper communication with external sites uses \textbf{HTTPS} only.
  \item Groq API requests are authenticated via \textbf{Bearer token} in the
        \texttt{Authorization} header.
\end{itemize}

% ─── CHAPTER 6: DATA REQUIREMENTS ────────────────────────────────────────────
\chapter{Data Requirements}

\section{Logical Data Model}
The system persists data in five relational tables. Their logical relationships
are described here; the physical schema diagram appears in the System Design
Document.

\begin{longtable}{L{3cm} L{10cm}}
  \toprule
  \textbf{Entity} & \textbf{Description} \\
  \midrule
  \endhead
  \textbf{Product} &
    Represents one canonical electronics product (e.g., ``Apple iPhone 15 128GB'').
    Identified by a unique slug. One product can have multiple listings and
    price history entries. \\[3pt]
  \textbf{ProductListing} &
    Represents the presence of a product on a specific retail platform
    (Amazon, Flipkart, or Vijay Sales) with current price, rating, stock
    status, and source URL. Unique per (productId, source) pair. \\[3pt]
  \textbf{PriceHistory} &
    An append-only record of price observations over time. One row per
    (product, source, scrape event). Used to generate price-trend charts. \\[3pt]
  \textbf{AIAnalysis} &
    Stores the Groq LLM output for a product: summary, recommendation,
    best-source name, and a confidence note. Refreshed every 24~hours. \\[3pt]
  \textbf{SearchLog} &
    Logs each user search query with a result count and timestamp.
    Used to populate the dashboard ``Recent Searches'' table. \\
  \bottomrule
\end{longtable}

\section{Data Retention and Integrity}
\begin{itemize}[noitemsep]
  \item \texttt{Product.slug} carries a \texttt{UNIQUE} constraint enforced
        by the database; Prisma upsert operations use it as the conflict key.
  \item \texttt{ProductListing} has a composite unique constraint on
        \texttt{(productId, source)} ensuring at most one listing per source
        per product.
  \item \texttt{PriceHistory} rows are never updated or deleted; they form
        an immutable audit trail.
  \item \texttt{AIAnalysis} has no unique constraint; the application uses
        \texttt{findFirst} followed by \texttt{update} or \texttt{create} to
        manage at most one active analysis per product.
  \item \texttt{SearchLog} is retained indefinitely for analytics; no
        PII is stored (no user IDs or IP addresses).
  \item All timestamps are stored in UTC and displayed in the user's local
        time zone via JavaScript's \texttt{Intl} API.
\end{itemize}

% ─── CHAPTER 7: CONSTRAINTS, ASSUMPTIONS, AND DEPENDENCIES ───────────────────
\chapter{Constraints, Assumptions, and Dependencies}

\section{Constraints}
\begin{enumerate}[noitemsep]
  \item \textbf{Amazon Bot Detection:} Amazon India returns HTTP 503 or 500
        for programmatic fetch requests, regardless of User-Agent spoofing.
        The Amazon scraper will silently return null in production and this
        is a known limitation. Fixing this would require using the Amazon
        Product Advertising API (requires seller account) or a rotating
        proxy service (out of scope for this project).
  \item \textbf{Flipkart Rate Limiting:} The Playwright scraper introduces a
        significant latency overhead (25--45 seconds per product) due to
        full browser initialisation and JavaScript rendering.
  \item \textbf{Vijay Sales Catalogue:} Vijay Sales does not stock all
        electronics models. Products unavailable in their catalogue return
        null and are excluded from comparisons.
  \item \textbf{LLM Token Limit:} The Groq API call is capped at 450 output
        tokens; very long product descriptions may be truncated.
  \item \textbf{Single-Region Scope:} All prices are in INR and reflect the
        Indian market only. Currency conversion is outside scope.
\end{enumerate}

\section{Assumptions}
\begin{enumerate}[noitemsep]
  \item The developer has a valid Groq API key with sufficient quota for the
        expected usage (~50 analysis requests per day).
  \item Docker and Docker Compose are installed on the deployment host.
  \item The host machine has at least 4~GB RAM to run PostgreSQL, Next.js,
        and a Playwright Chromium instance concurrently.
  \item The \texttt{.env.local} file is populated with \texttt{DATABASE\_URL},
        \texttt{GROQ\_API\_KEY}, and \texttt{SCRAPE\_SECRET} before running
        the application.
  \item Vijay Sales and Flipkart do not significantly change their DOM or
        API structure during the project evaluation period.
\end{enumerate}

\section{Dependencies}
\begin{longtable}{L{3.5cm} L{2cm} L{7.5cm}}
  \toprule
  \textbf{Dependency} & \textbf{Version} & \textbf{Purpose} \\
  \midrule
  \endhead
  Next.js & 16.2.4 & Full-stack React framework \\
  React & 19.2.4 & UI rendering library \\
  Prisma & 7.7.0 & Database ORM \\
  \texttt{@prisma/adapter-pg} & 7.7.0 & PostgreSQL driver adapter for Prisma 7 \\
  Playwright & 1.59.1 & Headless browser for Flipkart scraping \\
  \texttt{groq-sdk} & 1.1.2 & Groq LLM API client \\
  TailwindCSS & 4 & Utility-first CSS framework \\
  Recharts & 3.8.1 & Chart components for price history visualisation \\
  Radix UI & various & Accessible headless UI primitives \\
  Sonner & latest & Toast notification library \\
  next-themes & latest & Dark mode theme management \\
  Docker / Compose & 24+ / 2.x & Container runtime for PostgreSQL \\
  \bottomrule
\end{longtable}

\chapter*{Revision History}
\begin{longtable}{C{1.5cm} C{2cm} L{4cm} L{5cm}}
  \toprule
  \textbf{Version} & \textbf{Date} & \textbf{Author} & \textbf{Changes} \\
  \midrule
  \endhead
  1.0 & April 2026 & Adrij Bhadra &
    Initial release. Covers all 10 functional requirements,
    NFRs, external interfaces, data model, and constraints. \\
  \bottomrule
\end{longtable}

\end{document}
